% !PRINTED 618
\weberpage{618}
%<<CONT
zu setzen ist. $r$ ist daher eine \gesperrt{ganze Zahl} in $\Omega_m$, und ihre Norm ist $=+1$; folglich ist $r$ eine Einheit.

$\Omega_m$ ist ein Normalk\"orper, der durch die $\mu$ Substitutionen
\[ s_n = (r, r^n) \]
in sich selber \"ubergeht.

Diese $\mu$ Substitutionen $s_n$ bilden eine \gesperrt{Abel'sche} Gruppe $\mathfrak{N}$, welche die \gesperrt{Galois'sche} Gruppe des K\"orpers $\Omega_m$ ist. Sie ist isomorph mit der gleichfalls durch $\mathfrak{N}$ zu bezeichnenden Gruppe der nach dem Modul $m$ genommenen Zahlclassen $n$ (\S.\,16).

Die Function $f(x)$ l\"asst sich in die $\mu$ Factoren $x - r^n$ zerlegen, und wir setzen daher
\weq{4}{f(x) = \prod\limits^{n} (x - r^n).}

Setzen wir darin $x=1$, und beachten, dass [nach (3)] $f(1) = q$ ist, so folgt
\weq{5}{q = \prod\limits^{n} (1 - r^n).}

Die Zahlen
\weq{6}{\sigma_n = 1 - r^n, \quad \sigma = 1 - r}
sind aber ganze Zahlen des K\"orpers $\Omega_m$, und die Formel (5) zeigt zun\"achst, \gesperrt{dass $q$ im K\"orper $\Omega_m$ in $\mu$ Factoren $\sigma_n$ zerlegbar ist.}

Wir beweisen zun\"achst, dass die $\mu$ Zahlen $\sigma_n$ mit einander associirt sind. Bedeuten $n, n'$ zwei durch $q$ nicht theilbare Zahlen, so k\"onnen wir eine nat\"urliche Zahl $a$ so bestimmen, dass
\[ n' \equiv a n \ (\mathrm{mod}\ m) \]
wird. Dann ist aber
\[ \frac{\sigma_{n'}}{\sigma_n} = \frac{1 - r^{an}}{1 - r^n} = 1 + r^n + r^{2n} + \cdots + r^{(a-1)n} \]
eine ganze Zahl, also $\sigma_{n'}$ durch $\sigma_n$ theilbar. Da hierin $n$ mit $n'$ vertauscht werden kann, so ist auch $\sigma_n$ durch $\sigma_{n'}$ theilbar, also sind beide Zahlen associirt (\S.\,138). Wenn daher $\varepsilon$ eine Einheit bedeutet, so ist nach (5):
\weq{7}{q = \varepsilon \sigma^{\mu}, \quad N(\sigma) = q.}

Es ist also die nat\"urliche Primzahl $q$ associirt mit der $\mu^{\text{ten}}$ Potenz einer ganzen Zahl $\sigma$ im K\"orper $\Omega_m$. Diese Zahl $\sigma$ ist aber auch im K\"orper $\Omega_m$ noch Primzahl, und zwar vom ersten Grade. Denn hat $\sigma$ irgend einen Theiler $\sigma'$, so muss die Norm von $\sigma'$ ein Theiler der Norm von $\sigma$ sein, also, wenn $\sigma'$ keine
%CONT>>
